% Part: computability
% Chapter: computability-theory
% Section: k-1

\documentclass[../../../include/open-logic-section]{subfiles}

\begin{document}

\olfileid{cmp}{thy}{k1}
\olsection{مثال لقابلية الاختزال}

لننظر في تطبيق لـ\olref[ppr]{prop:reduce}.

\begin{prop}
\ollabel{prop:k1}
لتكن
\[
K_1 = \Setabs{e}{\cfind{e}(0) \fdefined}.
\]
عندئذ تكون $K_1$ قابلة للتعداد حسابيًا، ولكنها غير قابلة للحساب.
\end{prop}

\begin{proof}
لما كان $K_1=\Setabs{e}{\lexists[s][T(e,0,s)]}$، فإن $K_1$ قابلة
للتعداد حسابيًا بحسب \olref[eqc]{thm:exists-char}.

ولبيان أن $K_1$ غير قابلة للحساب، لنبين أن $K_0$ قابلة للاختزال
إليها.

\begin{explain}
هذا دقيق قليلًا، لأننا لا نستطيع باستعمال $K_1$ إلا طرح أسئلة عن
حسابات تبدأ بدخل معين هو $0$. افترض أن لك صديقًا ذكيًا يستطيع الإجابة
عن أسئلة من هذا النوع (ويسمى أصدقاء كهؤلاء «عُرّافين»). ثم افترض أن
شخصًا أتاك وسألك هل $\tuple{e,x}$ في $K_0$، أي هل تتوقف الآلة $e$ عند
الدخل $x$. ومن الأمور التي تستطيع فعلها بناء آلة أخرى $e_x$ تتجاهل،
عند \emph{أي} دخل، ذلك الدخل وتشغل بدلًا منه~$e$ عند الدخل $x$.
وعندئذ يكافئ بوضوح سؤال توقف الآلة $e$ عند الدخل $x$ سؤال توقف الآلة
$e_x$ عند الدخل $0$، أو أي دخل آخر. فتسأل صديقك هل تتوقف هذه الآلة
الجديدة $e_x$ عند الدخل $0$؛ ويوفر جواب صديقك عن السؤال المعدل جواب
السؤال الأصلي. وهذا يعطي الاختزال المطلوب لـ$K_0$ إلى~$K_1$.
\end{explain}

باستعمال الدالة الجزئية الشاملة القابلة للحساب، لتكن $f$ الدالة
الثلاثية الحجج المعرفة بـ
\[
f(x,y,z) \simeq \cfind{x}(y).
\]
لاحظ أن $f$ تتجاهل دخلها الثالث تمامًا. اختر فهرسًا $e$ بحيث
$f=\cfind{e}[3]$؛ فيكون لدينا
\[
\cfind{e}[3](x,y,z) \simeq \cfind{x}(y).
\]
بحسب مبرهنة $s$-$m$-$n$، توجد دالة $s(e,x,y)$ بحيث يكون، لكل $z$،
\begin{align*}
\cfind{s(e,x,y)}(z) & \simeq \cfind{e}[3](x,y,z) \\
& \simeq \cfind{x}(y).
\end{align*}

\begin{explain}
بعبارة الحجة غير الصورية أعلاه، $s(e,x,y)$ فهرس الآلة التي تتجاهل،
عند أي دخل $z$، ذلك الدخل وتحسب $\cfind{x}(y)$.
\end{explain}

وبوجه خاص، لدينا
\[
\cfind{s(e,x,y)}(0) \fdefined \quad \text{إذا وفقط إذا} \quad
\cfind{x}(y) \fdefined.
\]
وبعبارة أخرى، يكون $\tuple{x,y}\in K_0$ إذا وفقط إذا كان
$s(e,x,y)\in K_1$. ومن ثم تكون الدالة $g$ المعرفة بـ
\[
g(w) = s(e,(w)_0,(w)_1)
\]
اختزالًا لـ$K_0$ إلى~$K_1$.
\end{proof}

\end{document}
